\iffalse
%\section*{Appendix I — Extended Proofs and Derivations}

This appendix contains the extended mathematical proofs, derivations, and
formal arguments omitted from the main text for clarity. It includes full
operator algebra derivations, stability proofs, and drift-boundedness arguments
that establish the substrate as a lawful, antientropic cognitive geometry.
\fi
%\subsection*{I.1 Proofs Omitted from Main Text}

%\textbf{I.1.1 Curvature Preservation Under IRT}

Let


\[
T_{\mathrm{IRT}} : \mathcal{M} \rightarrow \mathcal{M}
\]


be a contraction-governed operator satisfying the Meaning Invariant.

Curvature preservation is expressed as bounded variance:


\[
|\kappa(T_{\mathrm{IRT}}(c)) - \kappa(c)| \le \epsilon_{\max}.
\]



Since $\kappa(c) \ge \kappa_{\min}$, we have the lower bound:


\[
\kappa(T_{\mathrm{IRT}}(c)) \ge \kappa_{\min} - \epsilon_{\max}.
\]



The substrate enforces the legality condition:


\[
\kappa(T_{\mathrm{IRT}}(c)) \ge \kappa_{\min},
\]


ensuring that only transitions satisfying the invariant are admitted into the
ledger. Thus, curvature preservation is enforced by legality rather than
algebraic necessity.

\textbf{I.1.2 Manifold Legality Under GBS}

Given:


\[
T_{\mathrm{GBS}}(c) \in \mathcal{M},
\]


and GBS defined as a contraction that cannot map to $\partial \mathcal{M}$, we
obtain:


\[
d(T_{\mathrm{GBS}}(c), \partial \mathcal{M}) > 0.
\]



Thus, GBS preserves manifold legality.

\subsection*{I.2 Extended Operator Algebra}

\textbf{I.2.1 Closure Under Composition}

Let $T_i, T_j \in \mathcal{O}$ be local contraction operators with constants
$k_i, k_j < 1$ relative to a context anchor $c^* \in \mathcal{M}$. We show:


\[
T_j \circ T_i \in \mathcal{O}.
\]



\textit{Proof.}

1. \textbf{Contraction property.}
Under the local contraction framework, distance is evaluated relative to the
active semantic context anchor $c^*$:


\[
d(T_j(T_i(c)), T_j(T_i(c^*)))
\le k_j \, d(T_i(c), T_i(c^*))
\le k_j k_i \, d(c, c^*).
\]


Since $k_j k_i < 1$, the composition remains a strict local contraction.

2. \textbf{Curvature variance.}
Unconstrained composition could allow variance accumulation up to
$2\epsilon_{\max}$ by triangle inequality. The substrate eliminates this drift
by enforcing \emph{sequential validation}: a composite transformation
$T_j \circ T_i$ is admitted into the ledger graph if and only if each
sequential edge


\[
e_n = (u_n \rightarrow v_n, T_n, c^*_n, \Delta_n)
\]


independently satisfies the machine-layer metadata constraint:


\[
\Delta \kappa_n = |\kappa(v_n) - \kappa(u_n)| \le \epsilon_{\max}.
\]


This renders the composite trajectory drift-bounded by constraint, ensuring
that $\mathcal{O}$ is closed under composition in the runtime sense.

3. \textbf{Identity preservation.}
IPU operators preserve $\mathcal{R}$; other operators do not modify it.

4. \textbf{Manifold legality.}
Each operator maps $\mathcal{M}$ to itself.

5. \textbf{Load stability.}


\[
|\Delta L_n| \le \delta_{\max}
\]


for each operator.

Thus, the operator algebra is closed under composition.

\textbf{I.2.2 Commutativity Conditions}

Operators commute when they preserve the same invariants:


\[
T_{\mathrm{IRT}} \circ T_{\mathrm{GBS}}
= T_{\mathrm{GBS}} \circ T_{\mathrm{IRT}}.
\]


Both preserve curvature within $\epsilon_{\max}$ and maintain manifold legality.

\subsection*{I.3 Stability Proofs}

\textbf{I.3.1 Stability Under Recursion}

Let $\gamma(t)$ be a recursive trajectory. Stability requires:


\[
\kappa(\gamma(t)) \ge \kappa_{\min},
\qquad
L(\gamma(t)) \le L_{\max}.
\]



Under contraction-governed operators:


\[
|\kappa(\gamma(t+1)) - \kappa(\gamma(t))| \le \epsilon_{\max},
\qquad
|L(\gamma(t+1)) - L(\gamma(t))| \le \delta_{\max}.
\]



Since only lawful transitions are admitted into the ledger, we obtain:


\[
\kappa(\gamma(t)) \ge \kappa_{\min},
\qquad
L(\gamma(t)) \le L_{\max}
\]


for all $t$ within the lawful recursion horizon.

\textbf{I.3.2 Entropy Boundedness}

Semantic entropy is defined as:


\[
H(t) = H_0 + \int_0^t \gamma(\tau) \, d\tau.
\]



Curvature invariants ensure:


\[
\kappa(t) \ge \kappa_{\min},
\qquad
\gamma(t) \le \gamma_{\max}.
\]



Thus:


\[
H(t) \le H_0 + t \gamma_{\max}
\quad \text{for lawful recursion horizons } t \le T_{\max}.
\]



The substrate prohibits unbounded recursion, preventing entropy divergence.

\subsection*{I.4 Drift Boundedness Proofs}

\textbf{I.4.1 Drift Metric Boundedness}

The drift score is:


\[
D(t) = w_1 \Delta_{\text{inv}}
      + w_2 \Delta_{\text{id}}
      + w_3 \Delta_{\text{geo}}
      + w_4 \Delta_{\text{load}}.
\]



Under lawful transitions:


\[
\Delta_{\text{inv}} \le \epsilon_{\max},
\qquad
\Delta_{\text{id}} = 0,
\qquad
d(\gamma(t), \partial\mathcal{M}) \ge \epsilon,
\qquad
\Delta_{\text{load}} \le \delta_{\max}.
\]



Thus:


\[
D(t) \le w_1 \epsilon_{\max} + w_4 \delta_{\max}.
\]



Drift is strictly bounded.

\textbf{I.4.2 Drift Explosion Under Probabilistic Systems}

Probabilistic recursion yields:


\[
\Delta_{\text{inv}} > \epsilon_{\max}, \quad
\Delta_{\text{id}} > 0, \quad
d(\gamma(t), \partial\mathcal{M}) < 0, \quad
\Delta_{\text{load}} > \delta_{\max}.
\]



Thus:


\[
D(t) \rightarrow \infty.
\]



This demonstrates that drift-boundedness is unique to substrate-native cognition.

\subsection*{I.5 Cross-Layer Constraint Propagation (CLCP) — Extended Proofs}

Cross-Layer Constraint Propagation ensures that invariant preservation holds
not only locally (within a section) but globally (across hierarchical layers).
This appendix provides the extended derivations establishing CLCP as a lawful,
drift-bounded global constraint.

\subsubsection*{I.5.1 Curvature Projection Bound}

For consecutive sections $\sigma_i$ and $\sigma_{i+1}$, define the curvature
projection:



\[
\kappa_{i \rightarrow i+1}
    = \left| \kappa(\sigma_{i+1}) - \kappa(\sigma_i) \right|.
\]



Local legality ensures:



\[
|\kappa(\sigma_{i+1}) - \kappa(\sigma_i)| \le \epsilon_{\max}.
\]



Global legality requires:



\[
\kappa_{i \rightarrow i+1} \le \kappa_{\max},
\]



where $\kappa_{\max}$ is the global curvature envelope defined in the substrate
geometry layer. Thus, curvature drift across layers is bounded.

\subsubsection*{I.5.2 Identity-Boundary Projection Bound}

Identity regions must remain stable across layers. Define:



\[
\partial I_{i \rightarrow i+1}
    = d\!\left( I(\sigma_i), I(\sigma_{i+1}) \right).
\]



Local IPU legality enforces:



\[
\text{region\_id}(\sigma_i) = \text{region\_id}(\sigma_{i+1}).
\]



Global CLCP legality strengthens this to:



\[
\partial I_{i \rightarrow i+1} \le \partial I_{\max},
\]



ensuring that identity boundaries do not drift across conceptual scales.

\subsubsection*{I.5.3 Manifold-Distance Projection Bound}

Let:



\[
d_M(i,i+1) = \mathrm{dist}_M(\sigma_i, \sigma_{i+1}).
\]



Local GBS legality enforces:



\[
d(\sigma_{i+1}, \partial\mathcal{M}) > 0.
\]



Global CLCP legality requires:



\[
d_M(i,i+1) \le d_{M,\max},
\]



ensuring that transitions remain within the global manifold envelope.

\subsubsection*{I.5.4 Load-Gradient Projection Bound}

Define:



\[
\nabla L_{i \rightarrow i+1}
    = \left| L(\sigma_{i+1}) - L(\sigma_i) \right|.
\]



Local LRP legality enforces:



\[
\nabla L_{i \rightarrow i+1} \le \delta_{\max}.
\]



Global CLCP legality requires:



\[
\nabla L_{i \rightarrow i+1} \le \nabla L_{\max},
\]



ensuring that load gradients remain stable across layers.

\subsubsection*{I.5.5 Global Drift-Boundedness}

Combining the four projection bounds yields:



\[
\kappa_{i \rightarrow i+1} \le \kappa_{\max}, \qquad
\partial I_{i \rightarrow i+1} \le \partial I_{\max}, \qquad
d_M(i,i+1) \le d_{M,\max}, \qquad
\nabla L_{i \rightarrow i+1} \le \nabla L_{\max}.
\]



Thus, CLCP enforces:



\[
\text{global drift} = 0,
\]



ensuring that the substrate remains scale-invariant and globally lawful.


\medskip

\noindent
\textit{This completes Appendix I.}
% \fi